\section{A Corrected Split}
\label{sec:corrected-split}

% Every number in this section traces to a committed run directory. Sources, by
% subsection:
%   4.1  runs/20260804_burst_aware_split_04/summary.md; tables/t2_split_properties
%   4.2  runs/20260801_device_split/summary.md (filename taxonomy);
%        runs/20260804_burst_aware_split_04/config.json (tau, seed, eps)
%   4.3  runs/20260804_burst_aware_split_04/{summary.md, groups_moved.csv,
%        moves.csv}; tables/t2_split_properties (39 of 61 -> 61 of 61)
%   4.4  runs/20260804_burst_aware_tau_sweep/summary.md
%   4.5  runs/20260803_corrected_split_02/summary.md (image-level allocator);
%        runs/20260804_burst_aware_split_03 via the sweep row at tau=30;
%        runs/20260814_figure_near_duplicate/per_box_centre_shift.csv
%   4.6  runs/20260804_build_corrected_dataset_02/summary.md
%   4.7  runs/20260804_duplicate_contamination/summary.md;
%        runs/20260805_consecutive_counter_pairs/
%   4.8  README.md licence section; data/kaggle/latency_history/pass2.json
%
% UNSUPPORTED BULLET, LEFT AS A TODO RATHER THAN WRITTEN. See 4.6 below: the
% drafted bullet claimed 364 same-camera cross-split pairs within 15 s in both
% the published and the corrected split. No committed run reports that figure.
% The nearest committed measurements are different quantities and must not be
% substituted for it:
%   - runs/20260803_corrected_split_02/ reports cross-split adjacent pairs
%     537 before / 600 after, but for the SUPERSEDED image-level allocator, not
%     the released one, and over all gaps rather than a 15 s window.
%   - runs/20260804_burst_aware_split_04/ reports 177 groups already straddling
%     the boundary before any move, which counts groups, not pairs.
% Producing the claimed figure needs a script and a run directory of its own.
%
% F3 LEGIBILITY -- MEASURED, NOT ESTIMATED. near_duplicate_pair is rendered at
% 11.0 x 7.4 in and its smallest text (the per-box shift table) is 6.9 pt at
% that size. Included at \figcol (3.5 in) that text renders at 2.2 pt; at
% \textwidth (6.5 in) at 4.1 pt. Both are below the 8 pt floor every other
% figure in this paper is built to. It is included below at \textwidth because
% that is the least bad option available, and it needs a re-render at target
% width before submission. Scaling cannot fix it.

The audit of Section~\ref{sec:dataset-audit} establishes a defect that is
repairable without collecting data: fifteen classes are unevaluable because the
sessions containing them were allocated wholly to training. This section
constructs a split that repairs it, states the constraints the repair had to
satisfy, and reports what it does and does not achieve. The construction is
released as a manifest rather than as a claim.

\subsection{What a correction has to satisfy}
\label{sec:correction-constraints}

Four constraints define the problem. Every class must reach at least five
instances in both the validation and the test partition, which is what makes it
evaluable at all. The split sizes must stay exactly 1478/438/205. No image may
be added, removed or modified. And test--train near-duplicate contamination,
which Section~\ref{sec:duplicate-contamination} measured at zero, must not
increase. The first is the objective; the other three are the price of a
correction that can be compared against the split it replaces.

The size constraint is the one that most needs defending, because it is the one
that costs the most. If the corrected training set were larger or smaller than
the published one, any later difference in accuracy between the two could be
attributed to the amount of training data rather than to its composition, and
the comparison in Section~\ref{sec:results} would answer a question nobody
asked. Freezing the three counts exactly means the only difference between the
two experiments is which images sit in which partition. It also makes the
correction harder: every image admitted into validation or test must be paid for
by another image returned to training.

\subsection{The grouping rule}
\label{sec:grouping-rule}

Whole capture bursts move as a unit; individual images never move alone.
Section~\ref{sec:duplicate-contamination} established that burst structure
exists in this dataset, and a burst is precisely a set of frames showing the
same scene seconds apart. Moving one frame out of a burst therefore places a
near-duplicate of a training image into an evaluation partition, which is the
contamination the fourth constraint forbids. Making the burst the atomic unit
of movement rules that out by construction rather than by detection: a group
that cannot be split cannot straddle the boundary.

Bursts are formed at a 15-second gap threshold and, critically, per camera
rather than globally. Sorting every image in the dataset by timestamp and
cutting at gaps would interleave frames from two photographers shooting at the
same moment on different devices, manufacturing groups that never existed as
capture events. Clustering within each device first and only then applying the
gap threshold keeps a group a record of one camera's continuous activity.

The grouping rule needs a taxonomy of filenames because the archive does not
carry one. Three families encode a timestamp and can be clustered in time; a
fourth cannot. Nine images carry an additional \texttt{\_NN} sub-index
appended to their timestamp, which resolves frames sharing a capture second.
The 189 untimestamped images admit no temporal clustering at all, so for those
the grouping falls back to scene content --- matching class inventory and box
geometry --- which groups by appearance rather than by time. That substitution
is weaker than a burst and is recorded as such: it is an inference about which
photographs show the same arrangement, not a measurement of when they were
taken.

\subsection{The allocation}
\label{sec:allocation}

The allocator begins from a per-class deficit: for each of the 61 classes, how
many further instances validation and test each need to reach five. Under the
published split 39 of the 61 classes already clear that bar in both partitions,
so 22 carry a deficit, and 15 of those 22 carry the maximum possible one, having
no instances in either partition to begin with.

Groups are then selected greedily, each step taking the group that closes the
largest number of outstanding deficits at once. This is why the repair is so
much cheaper than the deficit count suggests: a single capture group contains
several components photographed together, so admitting one group can lift
several classes over the bar simultaneously. Across the whole allocation 68
images move, and they repair all 22 classes --- taking the split from 39 of 61
to 61 of 61, as Table~\ref{tab:split-properties} records.

An equal number of images is returned to training, which is what holds the sizes
exactly. Only groups lying wholly inside one partition are eligible to be
returned, because returning part of a group would reintroduce the straddling
that Section~\ref{sec:grouping-rule} exists to prevent, and a group may be
returned only if removing it leaves every class still at or above five. The
dataset divides into 776 atomic groups, of which 177 already straddle a
partition boundary in the published split and are left untouched. Of the
remainder, validation admitted 18 groups and returned 10; test admitted 13 and
returned 6.

\subsection{Why 15 seconds}
\label{sec:tau-choice}

Class feasibility never bound the choice of threshold. At every value swept ---
15, 20, 25, 30, 35, 45 and 60 seconds --- all fifteen rescued classes retain two
or more qualifying groups, so the criterion the threshold was originally
introduced to satisfy is satisfied everywhere. Panel~(b) of
Figure~\ref{fig:tau-sweep} shows that row filled across the entire sweep, which
is the negative result that makes the rest of this subsection necessary: had
feasibility bound, the choice would have been made for us.

What bound it was the return budget. As the threshold grows, groups grow, and
the number of images the test partition can safely give back collapses --- 59 at
15 seconds, 26 at 25, 10 at 30, and 1 at 60 --- while the number it owes rises
in the opposite direction, from 15 to 31. The two cross between 25 and 30
seconds, and past that crossing the split cannot be rebalanced at all: at 30
seconds test owes 20 images and can return only 10, so the sizes end at
1458/438/225.

Fifteen seconds is the smallest value satisfying both constraints, and among the
three values that hold the sizes it moves the fewest images: 68, against 78 at
20 seconds and 80 at 25. It is not the value that moves the fewest images
overall. Thirty seconds moves 64, four fewer, and is excluded because it breaks
the sizes rather than because it is more expensive.

The criteria do not fail monotonically, and reporting only the chosen value
would hide that. Forty-five seconds is clean on contamination --- zero
test--train pairs at every epsilon under both scorings --- yet fails on sizes,
ending at 1448/438/235. Thirty-five and 60 seconds each carry one test--train
pair at $\epsilon = 0.05$ under aligned scoring, while both are clean under raw
scoring at the same threshold, which is why the scoring must be named whenever
that column is quoted. A threshold chosen by scanning upward until the first
failure would have stopped at the wrong place.

\begin{figure}[tp]
  \centering
  \includegraphics[width=\figcol,height=\figmaxh,keepaspectratio]{f4_tau_sweep.pdf}
  \caption{Why the released split uses $\tau = 15$~s. (a) the size constraint:
  at each $\tau$, the images the test split owes back after admitting whole
  groups, against the images it can safely return. The two cross between 25 and
  30~s, and past that point the split cannot be rebalanced --- at $\tau = 30$~s
  test owes 20 and can return only 10, ending at 1458/438/225 instead of
  1478/438/205. (b) the three admissibility criteria. Class feasibility is
  satisfied at every $\tau$ tested, so it never bound the choice; what bound it
  was the collapsing return budget in (a). $\tau = 15$~s is the smallest value
  meeting all three.}
  \label{fig:tau-sweep}
\end{figure}

\subsection{How the split was arrived at}
\label{sec:split-iterations}

The released split is the third of three attempts, and the two it replaced are
reported here rather than discarded. The sequence is evidence that the
constraints in Section~\ref{sec:correction-constraints} were binding on the
construction rather than assembled afterwards to describe whatever the final
allocator happened to produce.

The first attempt moved individual images. It succeeded on the objective and on
the sizes: 64 images moved, all 61 classes cleared the bar, and the split ended
at 1478/438/205 exactly. It failed on contamination, and the failure is
concrete. It moved \texttt{IMG\_20240220\_115315} into test while
\texttt{IMG\_20240220\_115316} --- captured one second later, showing the same
five components --- stayed in training. Across those five components the box
centres differ by at most 9.01 pixels on a 640-pixel canvas. That single
decision created 2 test--train near-duplicate pairs under raw scoring and 4
under aligned at $\epsilon = 0.05$, where the published split had none, which is
what established that an allocator blind to duplicates can satisfy every other
constraint and still make the split worse.

The second attempt made whole groups the unit of movement at a 30-second
threshold. It eliminated the contamination --- zero test--train pairs at every
epsilon under both scorings --- and broke the sizes instead, ending at
1458/438/225. Having each constraint satisfiable alone, and the two attempts
failing on different ones, is what motivated sweeping the threshold rather than
choosing it.

The third attempt swept the threshold across seven values and found that 15, 20
and 25 seconds all satisfy both constraints, of which 15 is the smallest and the
cheapest. That is the released split: $\tau = 15$~s, scene threshold
$\epsilon = 0.05$, seed 20260804.

Both superseded runs are retained in the repository with a banner recording why
they were replaced, and no number inside either was altered when they were
superseded. A run whose conclusion was overturned is still the only record of
what was actually computed at the time, and editing it would destroy the
evidence that the sequence above took place.

% F3 at \textwidth, not \figcol: see the legibility note in this file's header.
% This is a wide two-panel exhibit; even \textwidth leaves its shift table at
% 4.1 pt. It needs a re-render at target width before submission.
\begin{figure}[tp]
  \centering
  \includegraphics[width=\textwidth]{near_duplicate_pair.pdf}
  \caption{The pair that decided the grouping rule.
  \texttt{IMG\_20240220\_115315} and \texttt{IMG\_20240220\_115316}, captured
  one second apart, with their annotation boxes drawn and named. Across the five
  shared components the box centres differ by at most 9.01 pixels on the
  640-pixel canvas. The first allocator placed the left frame in test and left
  the right in training; the released split moves whole bursts, so both sit in
  training. The split shown under each panel is read from the built tree, not
  asserted, so the figure tracks whatever split is released.}
  \label{fig:near-duplicate-pair}
\end{figure}

\subsection{Verification}
\label{sec:split-verification}

The corrected dataset was materialised by copying from an unmodified v2 tree and
then verified by re-reading what had been written, under eight checks, all of
which pass. The image counts per partition are 1478/438/205; the copy loop's
intent matches what is on disk; no image lacks a label and no label is an
orphan; the total instance count is 12{,}937, unchanged; all 61 classes hold at
least five instances in both validation and test; every class identifier lies in
the range 0 to 60; and every copied file is byte-identical to its source.

The byte-identical check earns its place precisely because it is redundant with
the others under normal conditions. A truncated or partially written copy would
pass every count-based check in the list --- the file exists, it has a label, the
label parses, the class identifiers are in range --- while silently corrupting
the pixels a detector trains on. Comparing content hashes across all 2{,}121
images and their 2{,}121 label files is the only check in the set that would
catch it.

% TODO -- UNSUPPORTED, DO NOT WRITE UNTIL A RUN EXISTS. The drafted bullet here
% reported an independent check on the allocator: same-camera cross-split pairs
% within 15 s numbering 364 in the published split and 364 in the corrected one,
% showing the allocator neither created nor removed any. No committed run
% reports that figure, and the nearest committed numbers measure different
% things (see this file's header). The claim is worth making -- it is a check on
% temporal adjacency independent of the label-geometry contamination measure --
% but it needs a script and its own run directory first.

\begin{table}[t]
  \centering
  \caption{Properties of the published split and of the corrected split released here. Near-duplicate pair counts are at $\epsilon = 0.05$ with low-information pairs excluded, reported separately for raw and aligned scoring. The corrected split holds the image counts of the published one exactly, so no accuracy difference can be attributed to training-set size.}
  \label{tab:split-properties}
  \begin{tabular}{lrr}
    \toprule
    Property & Published & Corrected \\
    \midrule
    Images (train / val / test) & 1,478 / 438 / 205 & 1,478 / 438 / 205 \\
    Classes with $\geq$5 inst. in val and test & 39 of 61 & 61 of 61 \\
    Classes with zero val+test instances & 15 & 0 \\
    test$\leftrightarrow$train pairs (raw / aligned) & 0 / 0 & 0 / 0 \\
    val$\leftrightarrow$train pairs (raw / aligned) & 1 / 3 & 1 / 3 \\
    val$\leftrightarrow$test pairs (raw / aligned) & 0 / 0 & 0 / 0 \\
    Images reassigned & --- & 68 \\
    Grouping $\tau$ (s) & --- & 15 \\
    Seed & --- & 20260804 \\
    \bottomrule
  \end{tabular}
\end{table}


\subsection{What the correction does not fix}
\label{sec:correction-limits}

Validation--train contamination is unchanged. At $\epsilon = 0.05$ the published
split and the corrected one both carry 1 such pair under raw scoring and 3 under
aligned, and they are the same pairs: pre-existing in the data as distributed,
created by neither allocator and removed by neither. The zero reported in
Section~\ref{sec:duplicate-contamination} and preserved here describes the
test--train relationship only. Test contamination biases the headline metric and
validation contamination biases model selection; they are different failures and
a single cross-split figure that merged them would report neither.

The rescued classes rest on very few images each. Five instances is the
threshold at which a class becomes measurable, not the sample size at which a
per-class average precision becomes worth quoting with confidence. The
correction changes those classes from unmeasurable to measurable and claims
nothing beyond that, and per-class results in Section~\ref{sec:results} should be
read with the instance count beside them.

The duplicate measurement's scope carries over unchanged from
Section~\ref{sec:duplicate-contamination}: only pairs sharing an identical class
inventory are ever compared. Among the untimestamped images, three
consecutive-numbered pairs fall outside that scope, one of them across train and
test. \texttt{IMG\_5243} carries six classes to \texttt{IMG\_5244}'s five, the
extra being an Arduino-Uno, and the two shared single-instance components sit
109.6 and 86.3 pixels apart on the 640-pixel canvas. Inspected visually, the
scene is rearranged rather than re-photographed. That pair was never scored, and
never scored is not the same as scored and cleared.

Partial overlap between images is pervasive and no split removes it. Every image
in a capture session shares a background and draws from the same pool of
components, so a training image resembling a test image in background and in
several components is the normal condition of this dataset rather than a defect
of any particular allocation. The contamination figures are bounded claims about
identically inventoried pairs, and they are not a guarantee that no training
image resembles a test image.

The threshold is established for one seed. The allocator is greedy, and a
different seed explores a different corner of the search space; the sweep in
Section~\ref{sec:tau-choice} answers the question for seed 20260804 and does not
prove that 15 seconds is the right threshold in general. What it establishes is
that a threshold satisfying every constraint exists and that this one does.

\subsection{What is released}
\label{sec:split-release}

The split is released as a manifest assigning each image name to a partition.
Applied to a copy of ElectroCom61 v2 obtained from the Mendeley archive, it
reproduces the corrected split exactly, because the seed is fixed and the
allocation deterministic. The manifest is offered under CC BY 4.0, the licence
of the dataset it derives from, and the attribution to the dataset authors
travels with it \citep{electrocom61}.

A copy of the re-split images is also published alongside the manifest, as a
convenience for anyone who would otherwise download the archive and rebuild the
tree locally. That copy is a redistribution of the dataset's images, carries the
same CC BY 4.0 terms and the same attribution, and modifies nothing: the
corrected split changes which directory an image sits in and no image's pixels
or annotations. Both artefacts describe the same assignment, and the manifest is
the authoritative one --- it is what the verification in
Section~\ref{sec:split-verification} was run against.
